\begin{longtable}{|p{0.168\textwidth}|p{0.20\textwidth}|p{0.46\textwidth}|p{0.092\textwidth}|}
\caption{Coordinte, Dimension, and Variable attributes for GDS 2.0 netCDF data files}
\label{tab:variable-attributes} \\
\hline \endfirsthead \endhead \endfoot \hline \endlastfoot
\rowcolor{blue!25} \textbf{Variable Attribute Name} & \textbf{Format} & \textbf{Description} & \textbf{Source} \\ \hline
\rowcolor{violet!25} 
axis & String & For use with coordinate variables only. The attribute 'axis' may be attached to a coordinate variable and given one of the values “X”, “Y”, “Z”, or “T”, which stand for a longitude, latitude, vertical, or time axis respectively. See: \url{http://cfpcmdi.llnl.gov/documents/cfconventions/1.4/cfconventions.html#coordinate-types} & CF \\ \hline

\rowcolor{violet!25} 
bounds & string & An attribute for coordinate variables to describe the vertices of the cell boundaries and thus the intervals between cells.  Examples: ‘time:bounds = “time\_bnds”’, ‘XC:bounds = “XC\_bnds”’, and ‘YC:bounds = “YC\_bnds”’, with the later definitions: ‘time\_bnds:long\_name = “time bounds of averaging period”’, ‘XC\_bnds:long\_name = “longitudes of tracer grid cell corners”’, and ‘XC\_bnds:long\_name = “longitudes of tracer grid cell corners”’. & CF \\ \hline

\rowcolor{violet!25} 
c\_grid\_axis\_shift & string & Logical location of a variable relative to a tracer grid cell center in terms of fractional distance to the adjacent cell center along a given axis (x, y, or z). In the MITgcm's Arakawa C-grid, variables at 'u', 'v', and 'w' points are offset by +/- 0.5 relative to the tracer grid cell center. & xmitgcm \\ \hline

\rowcolor{violet!25} 
comment & string & Miscellaneous information about the variable or the methods used to produce it. & CF \\ \hline

\rowcolor{violet!25} 
coverage\_content\_type & string & An ISO 19115-1 code to indicate the source of the data, MD\_CoverageContentTypeCode [116]. Examples: “image,” “thematicClassification,” “physicalMeasurement,” “auxiliaryInformation,” “qualityInformation,” “referenceInformation,” “modelResult,” “coordinate”. & ACDD \\ \hline

\rowcolor{violet!25} 
long\_name & string & A free-text descriptive variable name. & CF, ACDD \\ \hline

\rowcolor{violet!25} 
positive & String & For use with a vertical coordinate variables only. May have the value “up” or “down”. For example, if an oceanographic netCDF file encodes the depth of the surface as 0 and the depth of 1000 meters as 1000 then the axis would set positive to “down”. If a depth of 1000 meters was encoded as -1000, then positive would be set to “up”. See the section on vertical-coordinate in [AD-3] & CF \\ \hline

\rowcolor{violet!25} 
standard\_name & string & Where defined, a standard and unique description of a physical quantity. For the complete list of standard name strings, see [AD-8]. \hspace{0.5cm}extbf\{Do not\} include this attribute if no \hspace{0.5cm}exttt\{standard\_name\} exists. & CF, ACDD \\ \hline

\rowcolor{violet!25} 
units & string & Text description of the units, preferably S.I., and must be compatible with the Unidata UDUNITS-2 package [AD-5]. For a given variable (e.g. wind speed), these must be the same for each dataset. Required for the majority of variables except mask, quality\_level, and l2p\_flags. & CF, ACDD \\ \hline

\rowcolor{violet!25} 
valid\_max & Expressed in same data type as variable & Maximum valid value for this variable once they are packed (in storage type). The fill value should be outside this valid range. Note that some netCDF readers are unable to cope with signed bytes and may, in these cases, report valid min as 127. Required for all variables except variable time. & CF \\ \hline

\rowcolor{violet!25} 
valid\_min & Expressed in same data type as variable & Minimum valid value for this variable once they are packed (in storage type). The fill value should be outside this valid range. Note that some netCDF readers are unable to cope with signed bytes and may, in these cases, report valid min as 129. Some cases as unsigned bytes 0 to 255. Values outside of \hspace{0.5cm}exttt\{valid\_min\} and \hspace{0.5cm}exttt\{valid\_max\} will be treated as missing values. Required for all variables except variable time. & CF \\ \hline

\end{longtable}
